\section*{\centering \KhOSml សេចក្តីសង្ខេប}
\vspace{2.5mm}
\addcontentsline{toc}{section}{\KhOSml សេចក្តីសង្ខេប}
{\KhOS
\hspace{1.27cm}
    % TODO: Write your Khmer abstract here.
    % Recommended: Use a Khmer word processor (e.g., LibreOffice with Khmer OS fonts)
    % to draft the text, then paste it here.
    %
    % The abstract should cover:
    % - Background and motivation of the Delta Robot research
    % - Methods used (mechanical design, kinematics, control)
    % - Key results and findings
    % - Contribution to robotics in Cambodia
    %
    % Example placeholder (replace with your actual Khmer text):
និក្ខេបបទនេះបង្ហាញពីការរចនា ការអភិវឌ្ឍ និងការផ្ទៀងផ្ទាត់ពិសោធន៍នៃប្រព័ន្ធ Delta robot លើកដាក់វត្ថុដែលប្រើប្រព័ន្ធចក្ខុវិស័យ ដោយប្រើ ROS 2 ក្នុងគោលបំណងផ្តល់នូវដំណោះស្រាយរ៉ូបូទិចដែលមានតម្លៃសមរម្យ និងមានប្រសិទ្ធភាព សម្រាប់ការអប់រំ និងការប្រើប្រាស់ឧស្សាហកម្មនៅប្រទេសកំពុងអភិវឌ្ឍន៍ ដោយផ្តោតជាពិសេសទៅលើប្រទេសកម្ពុជា។ Delta robot ដែលល្បីល្បាញដោយសាររចនាសម្ព័ន្ធចលនាស្របតទួល (parallel kinematic) និងសមត្ថភាពលើកដាក់វត្ថុដោយល្បឿនខ្ពស់ ត្រូវបានអនុវត្តជាប្រព័ន្ធស្វយ័តពេញលេញ ដែលបូករួមចក្ខុវិស័យកុំព្យូទ័រ ការគ្រោងផែនការចលនាក្នុងពេលវេលាជាក់ស្តែង និងការដំណើរការតាមបណ្តាញ CAN bus។
    \par
    \hspace{1.27cm}
    % [Add more Khmer paragraphs here]
    រ៉ូបូនេះត្រូវបានសាងសង់ដោយប្រើប្រដាប់ដំណើរការ RobStride RS-00 quasi-direct-drive ចំនួនបី ដែលត្រូវបានគ្រប់គ្រងតាម CAN bus ដំណើរការ ១ Mbit/s ស៊ុមរចនាសម្ព័ន្ធផ្លាស់ប្ដូររ (aluminum) និងបោះពុម្ព 3D ដោយមានប៉ារ៉ាម៉ែត្រធ្យូមគ្នា f = 157 mm, e = 35 mm, rf = 200 mm និង re = 400 mm ព្រមទាំង gripper ដ្រេសខ្យល់ (pneumatic solenoid gripper)។ ក្របខ័ណ្ឌកម្មវិធីប្រភពបើកចំហ រួមមាន Robot Operating System 2 (ROS 2) និង Python ត្រូវបានប្រើសម្រាប់ការគ្រោងផែនការចលនា ការគណនា kinematics និងការគ្រប់គ្រងពេលវេលាជាក់ស្តែង។ ប្រព័ន្ធចក្ខុវិស័យប្រើប្រាស់កាមេរ៉ា Intel RealSense D455 stereo depth ដែលផ្តល់ស្ទ្រីម RGB និង depth តម្រឹម (aligned) ដល់ ៣០ fps រួមជាមួយការរកឃើញវត្ថុក្នុងទំហំពណ៌ HSV (ដែលអាចប្ដូរជំនួស YOLO11n backbone បាន) សម្រាប់ការកំណត់ទីតាំងវត្ថុក្នុងពេលជាក់ស្តែង។ កូអរដោនេ pixel ដែលបានរកឃើញ ត្រូវបានព្យាករបញ្ច្រាសទៅជាកូអរដោនេ 3D ក្នុងស៊ុមកាមេរ៉ា ដោយប្រើម៉ូដែល pinhole ហើយបំប្លែងទៅស៊ុមមូលដ្ឋានរ៉ូបូតាមរយៈការបំប្លែង rigid-body ដែលបានកែតម្រូវ (calibrated)។
    \par
    \hspace{1.27cm}
    % [Add more Khmer paragraphs here]
    រ៉ូបូដែលបានអភិវឌ្ឍមានសេរីភាពចំណាកក្នុង ៣ ទិសដែលប្រតិបត្តិការក្រោម preset ចលនាមធ្យម ដោយ end-effector មានល្បឿនអតិបរមា vmax = 2000 mm/s និងសំទុះ amax = 5000 mm/s²។ គំរូ kinematic រួមមាន forward និង inverse kinematics ត្រូវបានទាញចេញតាមរបៀប analytic ដោយអនុវត្តតាមរូបមន្ត Trossen Robotics ហើយបានផ្ទៀងផ្ទាត់តាមការសាកល្បងរូបវ័ន្ត ដោយសម្រេចបាននូវកំហុស FK position ទាបជា ០,៥ mm ជាប់លាប់។ ការកែតម្រូវវ loop feedforward ដោយប្រើចក្ខុវិស័យ end-effector តាមដាន EE laser marker ក្នុងស៊ុមកាមេរ៉ា ហើយកែតម្រូវទីតាំងចូលទៅដំណើរការ (approach position) បន្តបន្ទាប់មុនពេលចុះ ដើម្បីបង្កើនភាពត្រឹមត្រូវក្នុងការលើករបស់នៅក្នុងទំហំការងារដែលមិនស្មើ (non-linear workspace)។ ការព្យាករខ្សែក្រវ៉ាត់ (conveyor belt prediction) ត្រូវបានអនុវត្តដោយប្រើទម្រង់ចលនាត្រីកោណ (triangular motion profile) ដើម្បីទូទាត់ចំពោះពេលវេលាធ្វើដំណើររ៉ូបូ នៅពេលលើកវត្ថុដែលកំពុងផ្លាស់ទី ក្នុងល្បឿនខ្សែក្រវ៉ាត់ ១២,៨៥ mm/s ។
    \par
    \hspace{1.27cm}
    សរុបមក ការស្រាវជ្រាវនេះបង្ហាញពីបំពង់ (pipeline) ពេញលេញពី vision ទៅ motion សម្រាប់ការលើកដាក់វត្ថុដោយស្វ័យ័ត ផ្អល់ platform Delta robot ដែលអាចចម្លងបាន (replicable) ដែលមានការចំណាយទាប និងលើកទឹកចិត្តការកសាងសមត្ថភាពក្នុងស្រុក ក្នុងវិស័យ robotics និង automation engineering នៅប្រទេសកម្ពុជា។
    % [Add concluding Khmer paragraph here]
}
